% 11-transparency.tex: watching the watchers.
\section{Watching the watchers: logs, witnesses and the split view}
\label{sec:transparency}

\subsection{The gap after anchoring}

Polaris periodically commits a batch of audit rows to a Merkle root and records the root in an
append-only table, so that an outside verifier can prove that the history under a root has not been
revised since. That leaves one gap. A single root proves the rows under it; it does not prove that
the operator did not later discard a root and publish a different one in its place. A transparency
log closes the gap by committing to the whole sequence of roots and letting anyone prove that the
sequence only ever grew. \sImpl

\subsection{The log}

The log's entries are the anchor roots in append order; nothing new is stored, because the
underlying table is already append-only. Over those entries the application builds a Merkle tree in
the style of Certificate Transparency (RFC 6962), with SHA3-256 and distinct prefixes for leaves and
interior nodes so that one can never be presented as the other. It publishes a signed tree head,
the log's commitment to its entire history at a size; a consistency proof that the size-$m$ tree is
a prefix of the size-$n$ tree, which exists only if the first $m$ entries were never reordered,
rewritten or dropped; an inclusion proof for any entry; and the entries themselves, so that a
monitor can replicate the log and recompute everything. The same machinery publishes two more logs:
the hashes of minted exchange receipts, and the hashes of anchored timestamps. There is no personal
data in any of it. \sImpl

\subsection{The monitor, and what it cannot see}

A monitor is a standalone program anyone runs. It keeps one piece of state, the last head it
accepted; each cycle it fetches the current head, verifies the signature against a key it was given
out of band, fetches a consistency proof from its last head to the new one, and either records the
new head or alerts on a rewritten entry, a fork, a shrunken tree, a head signed by the wrong key, or
a timestamp that went backwards. What it proves is consistency: the operator cannot rewrite history
without every monitor that saw the old head detecting it. What it cannot prove on its own is that
two monitors were shown the same head at the same size. A log can show one observer a private fork
of history. \sImpl

\figfile{transparency}{The log, the monitor, the witnesses, the gossip pool and the equivocation
proof. The log publishes a signed head and proofs. A monitor proves the log only appended. Witnesses
cosign heads they found consistent, a relying party requires a threshold of them on the exact head,
and witnesses gossip the heads they were shown; two log-signed heads of the same size with different
roots are a non-repudiable proof that the log lied. Heads can also be published into an external
append-only ledger that returns a receipt.}{fig:transparency}

\subsection{Witnesses, gossip and the equivocation proof}

A witness is a monitor that does two more things. It cosigns a head with its own key, only when the
head is an append-only extension of the last one it cosigned, and it gossips the log-signed head it
saw into a shared pool. Two uses follow, both verified offline. A witnessed checkpoint is a head
carrying cosignatures from at least $K$ distinct trusted witnesses, so a split view needs $K$
independent witnesses to equivocate, not just the log. An equivocation proof is two heads for one
log, both validly signed by the log's key, at the same size with different roots; the log signed
both, so it cannot repudiate the lie, and two gossiping witnesses produce the proof the moment they
compare notes. A witness that is itself shown a second head at a size it already cosigned refuses
and writes the proof. \sImpl

An external ledger completes the picture: the log publishes each head into an independent
append-only ledger and receives a receipt, the ledger's own signed head plus an inclusion proof, so
the log cannot use a head it has not publicly committed and the ledger cannot later drop it. A
ledger is itself an append-only log, so this reuses the same machinery and is monitored the same
way. The file-backed ledger is implemented and tested; the chain-backed drivers are declared and wait
on their interfaces. \sImpl{} for the file backend; \sFut{} for a chain.

\subsection{Who runs the witnesses}

The strength of this defence in practice is exactly the independence of whoever runs the
witnesses. What is built is the protocol: the cosignature format, the threshold rule, the
equivocation proof, and a witness daemon anyone can run, all drilled under real signatures every
release with three witnesses catching a fork two ways. What cannot exist until deployment is a set
of institutionally independent witnesses, a bank, a civil-society group, another authority,
actually running it. Today the witnesses are three processes in a drill. The repository states that
this is not a software deficiency but the absence of a deployment, and this paper repeats it,
because a transparency layer with witnesses run by the party it watches is not one. \sExt

\begin{layercard}{Layer card: transparency}
\qa{What it is}{Three RFC-6962-style logs over append-only tables; signed heads; consistency and inclusion proofs; a monitor; witnesses with cosignatures, a threshold and gossip; the equivocation proof; external-ledger publication with receipts.}
\qa{Why it exists}{A signed object proves what was said; only a watched sequence proves that what was said was not later unsaid, and only compared notes prove it was said the same way to everyone.}
\qa{What it trusts}{The log's key for the head; the witnesses' keys and their independence, which is a deployment fact.}
\qa{What it refuses}{A head with no consistency proof from the last one; a head not cosigned by the threshold; a receipt the ledger did not record; a second head at a size already seen.}
\qa{What it can see}{Roots and hashes. Never a person, never a payload.}
\qa{What it cannot see}{On its own, a split view: that is what the witnesses are for.}
\qa{If it fails}{A rewrite fails every monitor; a fork is a non-repudiable proof; a dishonest witness is one of $K$; a dropped published head fails the ledger's own consistency.}
\qa{Checked by}{\code{check\_transparency\_log}, \code{check\_transparency\_gossip}, \code{check\_transparency\_publication}, \code{check\_timestamp\_transparency}, the canonical-equivalence oracle for the signed head.}
\qa{Now or later}{Now, as a protocol with a drill. Independent witnesses are a deployment, later.}
\qa{Inside and outside}{Inside: federation. Outside: the institutions that exchange, sign and timestamp under this watch.}
\end{layercard}

\nextlayer{With authority, status and history checkable by strangers, institutions can build on the
credential: exchange requests, sign documents, attach time. The danger is that an exchange fabric
becomes the very message log the vocation forbids. The next layer is built as its inversion.}
